\section{Inspiration der Bibel}
\begin{otsblock}[Vorwort]
    Gestern haben im Bahnhof Brig eine Evangelisation durchgeführt. Dafür möchte ich unseren Jungen Leuten hier recht herzlich Danken. Ich muss ehrlich sagen, dass ich alleine nicht dem Mut aufgebracht hätte so etwas durchzuführen.

    Anton war letztes Jahr auf der gleichen Bibelschule wie ich im EBTC und dort haben wir den Kurs \enquote{Bibelstudium mit Gewinn} gemacht.\\
    Ich wusste nicht einmal, dass Anton auch im MNR war und so haben wir uns eigentlich nie unterhalten bis am Ende der Schule. Ich weiss nicht mehr wer, aber jemand hat mir gesagt, ich solle den Anton ansprechen, er würde mit einer Gruppe junger Leute regelmässig Evangelisationen durchführen.\\
    Am letzten Tag habe ich ihn dann angesprochen und das Resultat haben wir gestern gesehen. Es wurden über XXX Bibeln verteilt.

    Habe ich jetzt das alles von mir aus gemacht? Haben alle von sich auch aus gehandelt? Ich bin eher der Meinung, dass Gott das alles eingefädelt und alle Beteiligten  benutzt hat, um Schlussendlich eine Evangelisation in Brig durchzuführen.

    Kommen wir zum eigentlichen Thema die Bibel von Gott inspiriert.
\end{otsblock}
\begin{otsblock}[Einleitung]
    Die Bibel, obwohl sie von Menschen geschrieben wurde, enthält die Botschaft von Gott an uns Menschen und nicht die Botschaft von Menschen über Gott an uns Mitmenschen.\\
    In 66 Büchern hat Gott \betonung[b, p]{seine} Botschaft an uns von menschlichen Autoren aufschreiben lassen. In allen Einzelheiten führte Gott diese Männer, sodass sie genau das schrieben, was er ihnen zugedacht hatte. Das ist die Lehre der Inspiration. Natürlich sind jetzt nicht alle Theologen und Christlich der gleichen Meinung, darum gibt es die sogenannte \enquote{Theorie der Inspiration}.
    Ich will hier kurz sechs Theorien ansprechen.
\end{otsblock}
\subsection{Die wörtliche, vollständige Inspiration}
\subsection{Die Diktiermethode}
\subsection{Der inspirierte Gedanke}
\subsection{Teilweise Inspiration}
\subsection{Neoorthodoxe Sicht der Inspiration}
\subsection{Die naturalistische Inspiration}
\begin{otsblock}[Fazit]
    Es gibt viele Beweise, die die Inspiration der Bibel stützen, aber die Schrift selber behauptet das von sich. Ihre Kraft hat sich in dem veränderten Leben von Millionen Menschen gezeigt, die auf die Worte und Verheissungen der Schrift vertraut habe.

    Haben wir jetzt die Evangelisation gemacht, oder wurden wir von Gott geleitet und geführt? Hat Gott uns an die Strippe genommen und uns zum Bahnhof gezogen? Haben Schlussendlich nicht wir selber gehandelt? Wissen wir, wozu und wem dadurch geholfen wurde? Ich denke, die Schreiber der Bibel verstanden auch nicht immer, was sie da aufschrieben oder zu welchem Zweck es gut war, aber sie taten es.

    Eines Tages werden wir es Himmel erfahren.
\end{otsblock}